\begin{abcliste}
\abc The body is at the height of the tip, the shadow of the pointer on the vertical axis: it moves harmonically with the length of the pointer as amplitude, $\result{A = \AO}$, and with the angular velocity of the pointer.
\begin{align}
 \omega &= \omF = 2\pi\times\f = \om \approx \result{\omP{0}{3}} \\
 T &= \TF = \frac{1}{\f} = \result{\TP{0}{2}}
\end{align}

\abc The angle of the pointer from the right is $\omega t$, so the height of its tip is
\begin{align}
 y(t) &= A\,\sin(\omega t)
\end{align}
At $t = 0$ the body is at the centre, $y = 0$; during the first quarter turn the tip rises to the top. It is graph $\result{\text{(c)}}$: it starts at zero and goes up first, to \AO\ at $t = T/4 = \SI{0.125}{\second}$.

\abc
\begin{itemize}
\item (a) starts at zero but goes down first, $y = -A\,\sin(\omega t)$: the pointer turning the other way, clockwise.
\item (b) starts at the top, $y = A\,\cos(\omega t)$: the shadow on the horizontal axis (the horizontal position of the tip), not the height of the tip. Or: the pointer pointing up at $t = 0$.
\item (d) starts at the bottom, $y = -A\,\cos(\omega t)$: a wrong start, the pointer pointing down at $t = 0$.
\end{itemize}
\end{abcliste}
